\documentclass[11pt]{article}
\usepackage[a4paper,margin=1in]{geometry}
\usepackage{amsmath,amssymb,amsthm,booktabs}
\newtheorem{theorem}{Theorem}
\title{A signed-square descent and density-$99.9999999996\%$ theorem for Sharipov's second cuboid quintic}
\author{Computational research draft}
\date{August 4, 2026}
\begin{document}
\maketitle

\begin{abstract}
We prove that the constant coefficient of every hypothetical integral
quadratic factor is a signed square.  Enforcing its one global sign in exact
congruence certificates excludes a rational quadratic factor for a set of
primitive parameter pairs of natural density
$3917894203934863/3917894203948032$.  The result is unconditional but does not
resolve the infinite residual family, so Sharipov's conjecture remains open.
\end{abstract}

Let $Q_{p,q}(t)=H_{p,q}(t^2)$, where $H_{p,q}\in\mathbb Z[X]$ is the monic
quintic obtained from the standard second cuboid polynomial.

\begin{theorem}
If $X^2+AX+B$ divides $H_{p,q}(X)$ in $\mathbb Z[X]$, then
$B=\varepsilon d^2$ for some $d\in\mathbb Z$ and
$\varepsilon\in\{+1,-1\}$.
\end{theorem}

\begin{proof}
For a prime $\ell\mid pq$, put $k=v_\ell(pq)$.  Coprimality and direct
coefficient valuations give the lower Newton polygon with vertices
$(0,10k),(2,2k),(3,0),(5,0)$, also for $\ell=2,3$.  Thus the five root
valuations are $4k,4k,2k,0,0$.  The valuation of the product of any two roots
is even.  Since $B$ divides $-p^{10}q^{10}$, every prime valuation of $B$ is
even, proving the assertion.
\end{proof}

\begin{theorem}
Let $p,q$ range over positive coprime integers.  There is an explicit union of
projective congruence classes modulo all odd primes at most $137$, having natural
density
\[
\frac{3917894203934863}{3917894203948032},
\]
such that, on every class in the union,
$H_{p,q}$ has no quadratic factor over $\mathbb Q$.
\end{theorem}

\begin{proof}
If $H_{p,q}$ has a quadratic factor over $\mathbb Q$, monicity and Gauss's
lemma give a monic factor in $\mathbb Z[X]$.  Its reduction is a monic
quadratic divisor modulo every prime.  Weighted homogeneity shows that the
existence of such a divisor depends only on $(p:q)\in\mathbb P^1(\mathbb F_\ell)$.
Exhaustive exact division by all monic quadratics gives the complete class
lists in the accompanying certificate.  For illustration, some of the
sign-separated residual counts are below; ``both'' means that divisors of
both allowed signs occur in the same projective parameter class.

\begin{center}
\begin{tabular}{rrrrr}
\toprule
$\ell$ & $\ell+1$ & square & negative square & both\\
\midrule
43&44&20&12&4\\
67&68&40&20&12\\
71&72&8&24&4\\
83&84&24&24&8\\
\bottomrule
\end{tabular}
\end{center}

Projective classes are uniform for primitive pairs, and the finite local
conditions are independent by the Chinese remainder theorem.  Applying CRT
separately to each of the two global signs and then inclusion--exclusion gives
residual density
$13169/3917894203948032$, hence the stated complement.  Complete class lists and
independent exact verification algorithms accompany this draft.
\end{proof}

No finite set of reduction primes can remove every parameter class: if $M$ is
their product, the primitive pairs $p\equiv0\pmod M$, $q=1$ reduce to
$H_{p,q}(X)=X^3(X^2+6X+1)$ at every selected prime.  Thus the theorem is a
descent certificate, not a proof of the full conjecture.

\end{document}
